\chapter{The Mapping Pipeline}\label{ch:prototype}

The methods of Chapter~\ref{ch:methodology} are implemented as a five-stage
pipeline, released as an open repository so that every number in
Chapter~\ref{ch:results} can be regenerated from the raw standard documents.
This chapter describes the system's structure, what was built for the project
and what was reused, and the artefacts each stage produces.

\section{Architecture}

\begin{figure}[htbp]
\centering
\begin{tikzpicture}[
  node distance=5mm and 7mm,
  stage/.style={draw, rounded corners=2pt, minimum height=9mm, minimum width=24mm,
                align=center, font=\small},
  art/.style={font=\scriptsize\itshape, align=center},
  arr/.style={-{Stealth[length=2.4mm]}, thick}
]
\node[stage] (ex) {Extraction};
\node[stage, right=of ex] (gt) {Ground\\truth};
\node[stage, right=of gt] (map) {Mapping\\(7 methods)};
\node[stage, right=of map] (ev) {Evaluation};
\node[stage, right=of ev] (rep) {Reporting};
\draw[arr] (ex) -- (gt);
\draw[arr] (gt) -- (map);
\draw[arr] (map) -- (ev);
\draw[arr] (ev) -- (rep);
\node[art, below=2.5mm of ex] {provision\\JSON};
\node[art, below=2.5mm of gt] {gt.csv\\(107 pairs)};
\node[art, below=2.5mm of map] {predictions,\\similarity matrices};
\node[art, below=2.5mm of ev] {metrics, errors,\\confidence intervals};
\node[art, below=2.5mm of rep] {tables,\\figures};
\end{tikzpicture}
\caption{Pipeline stages and the artefact each stage produces. Every stage is
a separate command, so intermediate outputs can be inspected before the next
stage runs.}
\label{fig:pipeline}
\end{figure}

Figure~\ref{fig:pipeline} shows the five stages. The pipeline is a set of
Python modules run stage by stage from the command line; each stage reads the
files of the previous one and writes its results to a versioned data
directory. This staged design was a deliberate choice over a single
end-to-end script: extraction errors are cheaper to catch before mapping runs,
and the API-based methods (which cost money per call) are only invoked once
their inputs are verified.

\section{Stage descriptions}

\paragraph{Extraction.} Converts both standard PDFs to text with
\texttt{pdftotext} (Poppler), removes recurring ETSI page furniture, locates
provision identifiers with pattern matching, and writes one JSON record per
provision with identifier, section, title, text, and modality. Output: 69
records for \standardA{}, 72 for \standardB{}.

\paragraph{Ground truth.} Materialises the reviewed annotation of
Section~\ref{sec:gt} as \texttt{gt.csv}: source and target provision
identifiers, one of six relationship labels, a one-sentence justification,
and a confidence score. A companion module draws the stratified probability
sample, emits and ingests the annotation batches, and computes the agreement
statistics of Section~\ref{sec:gtvalidation}.

\paragraph{Mapping.} One module per method. The rule-based module computes
TF--IDF cosine similarity (scikit-learn) and security-keyword Jaccard scores.
The three local transformer modules share a common embedding helper that
loads the model (Hugging Face), embeds all provisions, normalises the
vectors, and computes the full $69 \times 72$ cosine similarity matrix. The
Gemini embedding module obtains vectors through the Google GenAI API; the
Gemini classification module builds one prompt per pair and submits them
through the batch API, then parses the returned labels. Every
similarity-based module stores both its thresholded predictions (CSV) and its
raw similarity matrix (NumPy), the latter enabling the threshold sweep of
Section~\ref{sec:thresholdresults} without re-running any model.

\paragraph{Evaluation.} Loads the ground truth and every prediction file,
computes detection and classification metrics
(Section~\ref{sec:evaluation}), bootstrap confidence intervals, per-method
confusion matrices, and threshold sweeps, and writes per-method error files
listing every missed pair and every wrong label with its provision
identifiers.

\paragraph{Reporting.} Renders the summary tables and the figures used in
Chapter~\ref{ch:results} (matplotlib), in print-ready PDF form.

\section{Built versus reused}

The extraction logic, the keyword baseline, the ground-truth tooling, the
evaluation harness with its bootstrap and sweep analyses, and the reporting
layer were developed for this project. The heavy lifting inside the mapping
methods deliberately reuses standard components: scikit-learn's TF--IDF
vectoriser, pretrained models from Hugging Face (MiniLM Sentence-BERT, BERT
base, SecureBERT), and Google's hosted Gemini models. No model was trained or
fine-tuned; assembling and evaluating existing components under controlled,
identical conditions is precisely the contribution.

\section{Reproducibility}

All randomised steps (sampling, bootstrap) use fixed seeds. A verification run
during the project reproduced every stored evaluation metric exactly;
individual embedding scores vary in the fourth decimal between runs because of
CPU thread scheduling in the tensor library, which never crosses a label
threshold. The API-based methods depend on hosted models and are archived
through their stored outputs instead. Appendix~\ref{app:reproduction} lists the exact commands, software
versions, and expected runtimes for a full reproduction.
